\scanpage{136}
\mspage{52}

En vertu du dual de loc cit 8.12.7, cela \deleted{implique}\inserted{\unsure{signifie en fait|signifie au fond}} que $f^{*}$ admette un adjoint à droite $f_{*}$

\[
f_{*}:M\unsure{\longrightarrow|\text{---}}N,
\tag*{(4.1.2)}
\]

lequel commutera donc aux petites $\underset{\longleftarrow}{\lim}$. La donnée de $f_{*}$ satisfaisant à cette propriété \unsure{détermine|déterminera} \unsure{donc|son} $f^{*}$ à \unsure{isom|iso} unique près. Mais en général, il ne suffit pas de se donner un tel $f_{*}$, pour pouvoir affirmer qu'il correspond à un $f$, i.e qu'il admette un adjoint à droite ; pour ceci, il faudrait qu'on sache que les foncteurs $M\unsure{\text{---}|\longrightarrow}\mathrm{Ens}$ qui commutent aux petites $\underset{\longleftarrow}{\lim}$ sont représentables, ce qui sera le cas si $M$ admet une petite sous-catégorie \uline{co}génératrice, i.e. si $M^{\mathrm{o}}$ est également une catégorie pseudo-topos.

Il est sans doute judicieux de procéder comme dans le cas des topos, et d'appeler \uline{morphisme de pseudo-topos} $f:\unsure{X|\mathcal X}\to\unsure{Y|\mathcal Y}$ un \unsure{quadruple|quadruplet}

\[
(f^{*},f_{*},\alpha,\beta)
\tag*{(4.1.3)}
\]
% archive sheet number: 135
